\label{sec:appendix_suggestions}
In this appendix suggestions to practitioners planning to carry similar studies are collected. These suggestions are based on the author's experience during the development of this thesis, shared hoping to facilitate future works in this field.

\subsubsection{Knowledge sharing sessions}
Elicitation sessions revealed fundamental alignment in expert judgment frameworks, with divergences emerging primarily from differences in thier specialized knowledge. This pattern was particularly pronounced in \gls{qi} assessments, where experts enjoyed greater interpretive freedom. Contextual awareness proved decisive: technological specialists often rated licensing statuses more favorably than colleagues attuned to country-specific political realities.

The \gls{smr} evaluations demonstrated how "insider" knowledge significantly influences outcomes. While most experts assigned similar baseline ratings, those with direct technology experience consistently applied more conservative assessments, reflecting awareness of recent, not well known, operational challenges. This underscores the value of multidisciplinary expert panels that incorporate both technical and contextual perspectives.

For these reasons, future studies should consider implementing knowledge-sharing sessions prior to individual elicitations. Such pre-elicitation exchanges would allow experts to present case-specific insights, potentially reducing deviations that stem from information asymmetries rather than genuine differences in professional judgment. This approach maintains individual assessment integrity while ensuring all experts operate from a common factual baseline.

\subsubsection{Methodological adherence in PILE-BWT elicitation}
During the \gls{pilebwt} elicitation, experts frequently defaulted to intuitive importance comparisons ("How much more important is A than B?") rather than the required compensatory trade-off formulation ("How much must A improve from its minimum to offset B's decline from minimum to maximum?"). Future applications must emphasize this distinction during expert briefings to maintain methodological rigor.

\subsubsection{Flexibility in value function correction}
During value function elicitation, when experts were asked to adjust the results of the "first elicitation", the rigid adherence to pre-defined indifference points from initial mid-splitting created unnecessary constraints. Future implementations should permit greater expert flexibility in adjusting these reference points during subsequent sessions.

\subsubsection{Possible improvements on QI approaches}
Regarding \glspl{qi} the use of a standardized reference alternative significantly improved consistency and is recommended for future studies. On the other hand, the "spectrum" approach for the "construction complexity" indicator (Section~\ref{sec:economic_indicators}) caused substantial expert confusion and should be replaced with conventional ranking methods.